% [原文 p.399]
\chapter{指数网络}

\ChapterAbstract{指数网络由 Eager、Selmani 与 Walcher 在 [122] 中引入，作为研究环面仿射 Calabi–Yau 三维流形（toric affine Calabi–Yau threefold）$X$ 上（带标架的）BPS B-膜的工具，它同时推广了 $4d$ $\mathcal{N}=2$ 规范理论中谱网络的构造，以及 [233] 中开创的 BPS 态的几何方法。由镜像对称，$X$ 上的 BPS B-膜（即 $X$ 上的凝聚层）对应于 $X$ 的镜像 $Y$ 中的特殊拉格朗日子流形。在局部镜像对称的框架下，这一对应进一步归结为对 $X$ 的镜像曲线 $\Sigma\subset(\C^\times)^2$ 中校准的拉格朗日圈的研究。选定一个投影 $\pi:\Sigma\to\C^\times$ 之后，人们便研究 $\C^\times$ 上的鞍点轨道（saddle trajectory），它们作为 $\Sigma$ 中校准圈在 $\pi$ 下的像出现。正是在这一设定中，指数网络出现了。

指数网络由 Banerjee、Longhi 与 Romo 在一系列论文 [38, 41, 39, 40] 中进一步发展；这一方向上更近的进展包括 [17, 101, 102, 169]。关于该主题最新的文献综述，我们还可参见 [42]。}

\section{定义与性质}

与谱网络类似，指数网络的形式定义依赖两类数据：
\begin{itemize}
\item 一部分几何数据，它依赖于由嵌入 $\Sigma\hookrightarrow(T^*C)^\times$ 诱导的黎曼曲面覆叠 $\Sigma\to C$，其中 $(T^*C)^\times$ 表示 $C$ 的余切丛去掉零截面。对每个相位 $\vartheta\in\R/2\pi\Z$ 的选取，这部分数据给出 $C$ 的一维子流形的一族校准。在大多数应用中取 $C=\C^\times$，这是局部镜像对称中自然出现的情形。
\item 一部分组合数据，它刻画由上述几何数据构造出的某个有效三维量子场论中的孤子含量。
\end{itemize}

% [原文 p.400]
对 $C=\C^\times_x$，有 $(T^*C)^\times\equiv\C^\times_x\times\C^\times_y$。该空间带有辛形式
\[
\omega=\frac{\dd x}{x}\wedge\frac{\dd y}{y}.
\]
考虑类 Liouville 的 1-形式
\[
\lambda=\log y\,\frac{\dd x}{x}=\log y\,\dd\log x;
\]
重要的是，它在 $(T^*C)^\times$ 上是多值的。可以定义由指数映射给出的无限循环覆叠 $\widetilde{\pi}:\widetilde{\Sigma}\to\Sigma$。于是我们有以下覆叠序列：
\[
\widetilde{\Sigma}\xrightarrow{\ \widetilde{\pi}\ }\Sigma\xrightarrow{\ \pi\ }\C^\times .
\]
对分歧覆叠 $\Sigma\to\C^\times_x$（其中 $\Sigma\subset(T^*C)^\times$）以及相位 $\vartheta\in\R/2\pi\Z$，$\C^\times_x$ 上的指数网络 $\mathcal{W}(\vartheta)$ 是一张由局部轨道构成的网，称为 E-壁（E-wall），它们是常微分方程
\[
\big(\log(y_j)-\log(y_k)+2\pi\ii n\big)\frac{\dd\log x(t)}{\dd t}\in \ee^{\ii\vartheta}\R_+
\]
的积分曲线，参数为 $t\in\R$。这些积分曲线由所有可能的 $(jk)_n$（包括 $j=k$）标号，并服从某些局部交叉规则，这些规则推广了谱网络的相应规则。从对数覆叠的观点看，这些规则是通过把多条开路径提升到接头之上的相继对数叶，并以尊重定向的方式匹配端点，从而形成一个闭合的圆而得到的。确切地说，指数网络中出现的接头（joint）或接点（junction）包括与 A 型谱网络类似的接头（涉及某些 $i,j,k,n$ 的 $(ij)_n$、$(ik)_n$ 与 $(jk)_n$ E-壁），还包括一类新的接头：$(ij)_n$ 与 $(ji)_m$ E-壁在其中相交，并产生无穷族新的 E-壁，如图 23.1 所示。由于这类新接头，指数网络不满足谱网络的局部有限性条件。

% [Tutor Note: 方程中的 $2\pi\ii n$ 来自 $\log y$ 在 $(\C^\times)^2$ 上的多值性：在对数覆叠的第 $n$ 层上，$\log y_k$ 与 $\log y_j$ 的差要加上 $2\pi\ii n$，这正是标号 $(jk)_n$ 的含义；把 $n$ 换成 $n+1$ 就是走到相邻的对数叶。核对 $j=k$ 的情形：此时 $\log y_j-\log y_k=0$，方程化为 $2\pi\ii n\,\dd\log x(t)/\dd t\in \ee^{\ii\vartheta}\R_+$；若按 $\R_+$ 为正实数的约定，则 $n=0$ 时无解，而 $n\ne0$ 时给出 $\dd\log x/\dd t$ 的辐角为常数，即 $x(t)$ 沿一条从原点出发的射线运动（$j=k$ 的 E-壁是直线）。这也解释了原文为何允许「包括 $j=k$」。]
\Fig[0.85\linewidth]{figures/fig_23_1.png}{23.1}{指数网络中出现的一类接头。}

与谱网络一样，对一般的 $\vartheta$，指数网络没有任何子网会提升为镜像曲线上的有限长闭合校准圈（从而是 $Y$ 中的紧拉格朗日对象）。

% [原文 p.401]
然而，有限网出现在非一般的 $\vartheta$ 值处；通过提升多条开路径并以尊重定向的方式匹配这些提升的端点，可以在 $\Sigma$ 中得到一个闭合圆，它反过来对应于 $Y$ 中的一个紧特殊拉格朗日圈。

与谱网络相比，指数网络的另一特殊性在于 E-壁在穿孔附近出现一种新的渐近行为：对数螺线。这具有物理意义：它解释了非紧 D4-膜。

\subsection*{镜像对称解释与 BPS 计数}

对环面仿射 Calabi–Yau 三维流形 $X$，指数网络刻画 $S^1\times X$ 上 M 理论的 BPS 谱。仿射环面 Calabi–Yau 三维流形 $X$ 的镜像 $Y$ 由一个多项式方程定义：
\[
Y=\{H(x,y)=uv\}\subset\C^\times_x\times\C^\times_y\times\C_u\times\C_v,
\]
其中 $H(x,y)$ 是一个 Laurent 多项式，其牛顿多边形是 $X$ 的环面图（受标架选取的含混性影响）。Calabi–Yau 流形 $Y$ 带有全纯 3-形式
\[
\Omega=\frac{\dd x}{x}\wedge\frac{\dd y}{y}\wedge\frac{\dd u}{u}.
\]
固定覆叠 $\pi:\Sigma\to\C^\times_x$ 的一个局部平凡化，使各分支可以标为 $y_k$（$k=1,\dots,d$），其中 $d$ 是覆叠的次数。上面的对数 1-形式 $\lambda$ 由沿投影 $Y\to(\C^\times)^2_{x,y}$ 的纤维对 $\Omega$ 积分得到。由构造，指数网络允许对 $Y$ 上的 BPS A-膜（从而是 $X$ 上的 BPS B-膜）计数，这可以用 $X$ 的 Donaldson–Thomas 不变量来表述。

% [Tutor Note: 本段给出镜像两侧膜类型的对照，可作阅读时的核对清单：镜像对称把 B-型膜（凝聚层、全纯对象）与 A-型膜（拉格朗日对象）互换，所以「$Y$ 上的 BPS A-膜」正对应「$X$ 上的 BPS B-膜」；原文摘要里说「$X$ 上的 BPS B-膜对应 $Y$ 中的特殊拉格朗日子流形」，与本段说法一致，只是出发点不同。再往下：由局部镜像对称，$X$ 的（局部）BPS 态归约为镜面曲线 $\Sigma\subset(\C^\times)^2$ 上的校准圈——$\Sigma$ 是复 1 维的，故这是实 1 维圈；而 $Y$ 是三维 Calabi–Yau，其中的紧特殊拉格朗日圈是 3 维的，二者之间差的正是 $\Omega$ 沿 $uv=H$ 纤维方向积分所得的 $\lambda$（其相位即校准相位 $\vartheta$）。核对方法：把 ODE 里的 $\log y_j-\log y_k$ 换成 $\lambda_j-\lambda_k$ 即得谱网络的相位条件，可见 $2\pi\ii n$ 项确由 $\log$ 的多值性带出。]
指数网络的计数还容许一个规范论的解释，不过它出现在 $5d$ $\mathcal{N}=1$ 理论的语境中，而不是谱网络所在的 $4d$ $\mathcal{N}=2$ 理论。具体而言，$S^1\times\R^4$ 上的 $5d$ $\mathcal{N}=1$ 理论对应于 $\C^\times$ 上的指数网络。Seiberg–Witten 曲线 $\Sigma$（在此语境中有时称为 Seiberg–Witten–Nekrasov 曲线 [283]）由多项式方程
\[
\{F(x,y)=0\}\subset(\C^\times)^2=(T^*\C^\times)^\times
\]
定义，其中 $F$ 是 Laurent 多项式，并且它可以等同于某个环面仿射 Calabi–Yau 三维流形的镜像曲线。

在实践中，指数网络已对一些简单的环面 Calabi–Yau 三维流形被显式构造出来，包括锥流形（conifold）[122]、局部 $\PP^2$ [122, 39] 以及局部 Hirzebruch 曲面 $\mathbb{F}_0$ [41]。

% [原文 p.402]
\section{路径提升与（非）阿贝尔化}

与谱网络类似，指数网络为同伦不变的路径提升提供了一个框架。具体地说，$\C^\times$ 上从属于分歧覆叠 $\Sigma\to\C^\times$ 的任一指数网络都诱导一个映射，它把 $\C^\times$ 上的每个同伦类对应到 $\Sigma$ 上（或在满足某些约束的前提下对应到其覆叠 $\widetilde{\Sigma}$ 上）的一族同伦类 [38]。

这一路径提升映射支撑着一个形式的（非）阿贝尔化过程，与谱网络的情形密切平行。设 $\mathcal{L}$ 是 $\widetilde{\Sigma}$ 上带有平坦阿贝尔联络的线丛。把它前推到 $\C^\times$ 得到一个无限秩丛，带有形式非阿贝尔平坦联络，即一种平行移动的概念。投影 $(T^*\C^\times)^\times\to\C^\times$ 的纤维就是 $\C^\times$ 本身，这一事实表明：由指数网络定义的（非）阿贝尔化自然地与环路群的特征簇相联系，而不是与有限维李群的相应对象相联系。

% [Tutor Note: 原文说路径提升把 $\C^\times$ 上的一个同伦类对应到 $\Sigma$ 上的一「族」同伦类，这与谱网络不同：谱网络的提升在选定起点所在的叶后是单值的；这里对数覆叠 $\widetilde\Sigma\to\Sigma$ 的叶由 $\Z$ 编号（见覆叠序列 $\widetilde\Sigma\xrightarrow{\widetilde\pi}\Sigma\xrightarrow{\pi}\C^\times$），而穿过一条 $(jk)_n$ 壁会把叶指标平移 $n$，故同一个基类可产生无穷多个提升类。若要逐条核对，可先按图 23.1 的 $(ij)_n$、$(ji)_m$ 与新生成的 $(ij)_{n+(k-1)(n+m)}$ 等标号跟踪指标变化，再把一圈走完后总指标回到原处（闭合圆）。]
% [原文 p.403]
\chapter{三维网络}

\ChapterAbstract{在 [138] 中，Freed 与 Neitzke 为维数小于三的流形上的谱网络引入了抽象框架，并描述了三维 Chern–Simons 不变量的一个分层阿贝尔化过程。具体而言，他们证明：带主 $\SL_2(\C)$-丛与平坦联络的三维流形 $M$ 的 $\SL_2(\C)$-Chern–Simons 不变量，可以用 $M$ 的某个分歧二重覆叠的 spin $\C^*$-Chern–Simons 不变量来计算。此外，作者们把三维（$3d$）谱网络的构造推广到群 $\GL_2(\C)$ 与 $\PSL_2(\C)$。}

\section{Chern–Simons 场论}

Chern–Simons 理论是一种三维拓扑规范理论，给出 $3d$ 流形的拓扑不变量。它可以表述为 [35] 中描述的 $3d$ 拓扑量子场论（TQFT）。具体而言，指数化的 Chern–Simons 不变量定义了一个从某个 3-配边范畴到一维复向量空间构成的群胚 $\underline{\mathrm{Line}}_{\C}$ 的对称幺半函子。

\begin{remark}[24.1]
Yang–Mills 泛函
\[
S_{\mathrm{YM}}(A)=\int_M \tr F(A)\wedge\star F(A),
\]
其中 $M$ 是伪黎曼流形，$A$ 是 $M$ 上固定主 $G$-丛 $P$ 的联络，$F(A)$ 是 $A$ 的曲率，$\tr$ 表示 $G$ 的李代数上的一个 Killing 型，它通过 Hodge 星算子显式依赖于 $M$ 的度量。相反，对四维流形 $M$，泛函
\[
S_\Theta=\frac{1}{8\pi^2}\int_M \tr F(A)\wedge F(A)
\]
% [原文 p.404]
是拓扑的。这个泛函计算主丛 $P\to M$ 的一个特征类，并且值得注意的是，它不依赖联络 $A$。微分形式 $\tr F(A)\wedge F(A)$ 以 Chern–Simons 3-形式
\[
\tr\Big(A\wedge \dd A+\tfrac{2}{3}A\wedge A\wedge A\Big)
\]
作为局部原像。
\end{remark}

% [Tutor Note: 可据此核对原文两处 CS 3-形式的系数是否自洽。取矩阵值形式的约定 $F(\theta)=\dd\theta+\frac12[\theta,\theta]$（等价地 $\theta\wedge\theta=\frac12[\theta,\theta]$），则 $\theta\wedge F(\theta)=\theta\wedge\dd\theta+\theta\wedge\theta\wedge\theta$，而 $\theta\wedge[\theta\wedge\theta]=2\,\theta\wedge\theta\wedge\theta$，于是 $\theta\wedge F(\theta)-\frac16\theta\wedge[\theta\wedge\theta]=\theta\wedge\dd\theta+\big(1-\frac13\big)\theta\wedge\theta\wedge\theta=\theta\wedge\dd\theta+\frac23\theta\wedge\theta\wedge\theta$，正是紧接着给出的 Chern–Simons 3-形式。可见 $\frac16$ 与 $\frac23$ 这两个系数互为印证（$\frac16=\frac12-\frac13$）。]
设 $X$ 是三维流形，$P\to X$ 是主 $G$-丛，其中 $G$ 是只有有限多个连通分支的李群，其李代数 $\g$ 上带有一个 $G$-不变对称双线性形式 $\tr$（例如 $G$ 半单时的 Killing 型）。对 $P$ 上任一 1-形式 $\theta\in\Omega^1(P,\g)$，定义
\[
\eta(\theta)=\tr\big(\theta\wedge F(\theta)-\tfrac{1}{6}\theta\wedge[\theta\wedge\theta]\big)\in\Omega^3(P,\C),
\]
其中 $F(\theta)$ 是 $\theta$ 的曲率。3-形式 $\eta(\theta)$ 是闭的，可用它如下定义 $X$ 的一个不变量。设 $s:X\to P$ 是 $P$ 的一个截面，考虑积分 $\int_X s^*\eta(\theta)$。由于 $\eta(\theta)$ 闭，这个复数只依赖于截面的同伦类。然而，两个截面未必同伦，因此需要额外修改才能使所得的量不依赖 $s$ 的选取：对称双线性形式 $\tr$ 必须是整的，即属于 $H^4(BG,\Z)$，其中 $BG$ 表示 $G$ 的分类空间。此外，这个不变量只在模 $2\pi\ii$ 的意义下有良好定义。

\begin{definition}[24.2]
带主 $G$-丛 $P\to X$ 与 $\theta\in\Omega^1(P,\g)$ 的三维流形 $X$ 的指数化 Chern–Simons 不变量定义为
\[
\mathcal{F}_G(X,\theta)=\exp\Big(2\pi\ii\int_X s^*\eta(\theta)\Big)\in\C^\times,
\]
其中 $s$ 是 $P$ 的一个截面。指数化 Chern–Simons 不变量不依赖 $s$ 的选取。
\end{definition}

现在描述维数为 3 的 Chern–Simons TQFT。TQFT 的基本想法是提供「带边流形的表示论」。更确切地说，维数为 $d$ 的 TQFT 给 $(d-1)$ 维闭定向流形配上向量空间，并给 $d$ 维定向流形配上线性映射。该线性映射从入边（即定向与诱导定向相符的边界分支）所对应的向量空间，映入出边所对应的向量空间。这些 $d$ 维流形称为配边（bordism）。

这个对应只应依赖流形的拓扑，从而保证流形粘合时线性映射可以复合，并且必须是乘性的。这意味着，两个 $(d-1)$ 维流形的不交并所对应的向量空间，是各自向量空间的张量积。特别地，TQFT 给每个闭 $d$ 维流形配上一个标量拓扑不变量。

% [原文 p.405]
\Fig[0.6\linewidth]{figures/fig_24_1.png}{24.1}{TQFT 的示意图。}

Chern–Simons $3d$ TQFT 给带丛与平坦联络的闭二维流形配上复直线（即一维复向量空间），并配上对应于配边的这些复直线之间的同构。对带联络 $\theta$ 的闭三维流形 $X$，不变量由指数化 Chern–Simons 不变量 $\mathcal{F}_G(X,\theta)$ 给出。

这个 $3d$ TQFT 的精确定义相当复杂。借助 $3d$ 谱网络，Freed 与 Neitzke 能够利用 $X$ 的某个二重覆叠的 spin $\C^*$-Chern–Simons 不变量，来计算 $G=\SL_2(\C)$ 时的 Chern–Simons 不变量。

\section{三维谱网络及其应用}

Freed 与 Neitzke 所用的基本概念是底流形的 SN-分层（SN-stratification）。

\begin{definition}[{{{{{[138, 定义 4.15]}}}}}]
设 $X$ 是带角的流形，维数至多为三。$X$ 的一个 SN-分层由以下数据组成：
\begin{enumerate}
\item 一个余维 0 的光滑流形，
\item 一些余维 1 的子流形，
\item 两类余维 2 的子流形：$b$ 型（其中 $b$ 代表「branch」（分支）），其链环（linking manifold）是恰有 3 个标记点的圆周；以及 $a$ 型，其链环是带有任意多个标记点的圆周，
\item 两类余维 3 的子流形：$b$ 型，其链环是由四面体诱导出分层的 2-球面；以及 $a$ 型，其链环是由顶点均为 $a$ 型的三价图分层的 2-球面。
\end{enumerate}
\end{definition}

注意，区分 $a$ 型与 $b$ 型是数据的一部分：例如在余维 2 中，可以存在 $a$ 型的层，其链环圆周恰有 3 个标记点。不同的层记作 $M_{-1}$、$M_{-2a}$ 等等。定义 $M_{\ge -3a}:=M\setminus(M_{-2b}\sqcup M_{-3b})$，它表示 $M$ 中去掉分支轨迹（branch locus）的部分。

% [原文 p.406]
\begin{definition}[{{{{{[138, 定义 4.24 (i)]}}}}}]
设 $M$ 是带边紧流形，维数至多为 3，并带有 SN-分层，使得 $M\setminus\partial M=M_0\sqcup M_{-1}\sqcup M_{-2a}\sqcup M_{-3a}\sqcup M_{-2b}\sqcup M_{-3b}$。从属于 $M$ 上该分层的谱网络 $(\pi,s)$ 由以下部分组成：
\begin{itemize}
\item 一个二重覆叠 $\pi:\widetilde{M}_{\ge 3a}\to M_{\ge 3a}$，它在 $M_{-2b}$ 中每一点的链环上的限制都是非平凡的，
\item $\pi$ 在 $M_{-1}\sqcup M_{-2a}\sqcup M_{-3a}$ 上的一个截面 $s$。
\end{itemize}
\end{definition}

最后，我们可以定义把 Chern–Simons 不变量阿贝尔化所需的数据。设 $T$ 是 $\SL_2(\C)$ 中的极大环面，$H$ 是它的正规化子。

\begin{definition}[{{{{{[138, 定义 4.24 (ii)]}}}}}]
设 $M$ 如前一定义。我们把以下数据定义为分层阿贝尔化数据（stratified abelianization data）：
\begin{itemize}
\item 一个带平坦联络的主 $\SL_2(\C)$-丛 $P\to M$，
\item 一个带平坦联络的主 $H$-丛 $Q\to M_{\ge 3a}$，以及 $M_0$ 上的一个平坦同构 $Q\cong P$，
\item $M_{\ge -3a}$ 上的一个二重覆叠的同构 $\widetilde{M}_{-3a}\cong Q/T$。
\end{itemize}
\end{definition}

% [Tutor Note: 原文定义 24.4 与 24.5 中的 $M_{\ge 3a}$ 疑为 $M_{\ge -3a}$ 之误，理由是：（i）全书层用负余维标号——$M_{-1}$（余维 1）、$M_{-2a}$（余维 2 的 $a$ 型）、$M_{-3a}$（余维 3 的 $a$ 型）等，$M_0$ 是余维 0 层；（ii）$M_{\ge-3a}=M\setminus(M_{-2b}\sqcup M_{-3b})$ 恰是去掉分支轨迹后、二重覆叠可以有定义的那部分，而二重覆叠 $\pi$ 的定义域理应就是它；（iii）同一定义 24.5 的第三条自身写作 $\widetilde{M}_{-3a}\cong Q/T$ over $M_{\ge -3a}$，与 $M_{\ge 3a}$ 不一致。此处照原文译出。]
存在与（截角）四面体相伴的标准 SN-分层。因此，每个已三角剖分的三维流形都被赋予一个标准 SN-分层，它推广了（部分）出现在带全纯二次微分的黎曼曲面上的二维谱网络。Freed 与 Neitzke 的主要结果是一个 Chern–Simons 不变量的整体阿贝尔化过程，它表述为可逆场论之间的同构：

\begin{theorem}[{{{{{[138, 定理 5.106]}}}}}]
设 $M$ 是带分层阿贝尔化数据的闭定向三维流形。则 $M$ 上的 $\SL_2(\C)$-Chern–Simons 场论同构于 $\widetilde{M}$ 上的 $\C^*$-spin Chern–Simons 理论与一个 spin $\Z/2\Z$-Chern–Simons 理论的张量积，后者不依赖于 $M$ 上的丛。
\end{theorem}

作为一个应用，Freed 与 Neitzke 给出了 Chern–Simons TQFT 的详细描述。特别地，他们确定了与边界曲面相伴的复直线的典范同构，并对闭三维流形计算了指数化 Chern–Simons 不变量。

总之，三维谱网络提供了一种把 Chern–Simons 不变量阿贝尔化的方法，可能引向许多令人兴奋的发展。由该论文产生了若干公开问题，特别是以下几个：
\begin{itemize}
\item 推广 $3d$ 谱网络，使其涵盖 $2d$ 谱网络的完整范围。
\item 把该框架推广到高阶李群。
\item 把 $3d$ 谱网络与 BPS 态或 WKB 近似联系起来。
\end{itemize}
